\documentclass{exos}
\usepackage{main}

\begin{document}

\begin{pycode}
import string
def encrypt(s,decalage):
  n = len(string.ascii_letters)
  for char in s:
    index = string.ascii_letters.find(char)
    if index != -1:
      new_char = string.ascii_letters[(index + decalage) % n]
      new_char = new_char.capitalize()
      print(new_char, end="")
    else:
      print(char, end="")

s = r'''Bonjour ! Vous avez reussi a dechiffrer le code secret. Le texte est volontairement long pour pouvoir trouver les lettres frequentes. L'indice pour cette epreuve est quarante-sept.'''

d = 7
\end{pycode}

Le code secret ci-dessous a été codé grâce au \textit{chiffrement César}. Chaque lettre a été décalé d'un certain nombre $n$ de cran. Par exemple, si le nombre est $3$, alors $A \to D$, $B \to E$, \dots, $Y \to B$, $Z \to C$.

\textbf{Mais ici, on ne sait pas de combien de crans le message est décalé. Arriverez-vous à le déchiffrer ?}

\vspace*{1cm}

\pyc{encrypt(s,d)}

\vspace*{1cm}

(Indice : Le lettre la plus frequente dans les mots en français est le E. Quelle est la lettre la plus fréquente ici ?)

\vspace*{1cm}
\hrulefill
\vspace*{1cm}

Le code secret ci-dessous a été codé grâce au \textit{chiffrement César}. Chaque lettre a été décalé d'un certain nombre $n$ de cran. Par exemple, si le nombre est $3$, alors $A \to D$, $B \to E$, \dots, $Y \to B$, $Z \to C$.

\textbf{Mais ici, on ne sait pas de combien de crans le message est décalé. Arriverez-vous à le déchiffrer ?}

\vspace*{1cm}

\pyc{encrypt(s,d)}

\vspace*{1cm}

(Indice : Le lettre la plus frequente dans les mots en français est le E. Quelle est la lettre la plus fréquente ici ?)

\vspace*{1cm}
\hrulefill
\vspace*{1cm}

Le code secret ci-dessous a été codé grâce au \textit{chiffrement César}. Chaque lettre a été décalé d'un certain nombre $n$ de cran. Par exemple, si le nombre est $3$, alors $A \to D$, $B \to E$, \dots, $Y \to B$, $Z \to C$.

\textbf{Mais ici, on ne sait pas de combien de crans le message est décalé. Arriverez-vous à le déchiffrer ?}

\vspace*{1cm}

\pyc{encrypt(s,d)}

\vspace*{1cm}

(Indice : Le lettre la plus frequente dans les mots en français est le E. Quelle est la lettre la plus fréquente ici ?)

\end{document}